\documentclass{amsart}
\usepackage{amsrefs}
\usepackage[T1]{fontenc}
\usepackage{graphicx} % Required for inserting images
\usepackage{amsfonts, amsmath, amssymb}
\usepackage[english]{babel}
\usepackage{framed}
\usepackage{array}
\numberwithin{equation}{section}
\newtheorem{theorem}{Theorem}
\newtheorem{lemma}{Lemma}
\newtheorem{conj}{Conjecture}
\newtheorem{coro}{Corollary}
\newcommand{\eps}{\varepsilon}
\newcommand{\Z}{\mathbb Z}
\newcommand{\N}{\mathbb N}
\newcommand{\R}{\mathbb R}
\newcommand{\C}{\mathbb C}
\newcommand{\Q}{\mathbb Q}
\def\legendre#1#2{\genfrac(){0.2pt}0{#1}{#2}}
\newcommand{\minilegendre}[2]{\genfrac{(}{)}{}{}{#1}{#2}}
\title{On series involving $\binom{4k}{k}^{-1}$}
\author{Roman Le Lan}
\email{roman.lelan1@gmail.com}
\begin{document}
\begin{abstract}
In this article, we prove two conjectures of Sun concerning the evaluation of two series involving the reciprocal of the binomial coefficient $\binom{4k}{k}$. The values of these series are related to $\pi$, $\log3$, and the Gieseking constant $\text{Cl}_2(\pi/3)$. Our proofs rely on the Beta function and the Clausen function.
\end{abstract}
\maketitle

\section{Introduction}
Recently, Sun published a paper \cite{series2} in which he mainly considers series of the form
\[\sum_{k=1}^\infty \frac{ak^2+bk+c}{k(3k-1)(3k-2)m^k\binom{4k}{k}}\]
where $a, b, c \in \Z$ and $m$ is a nonzero rational number.
Many evaluations of such series have been established. For example, Au \cite{Au} showed that
\begin{equation}
\sum_{k=1}^\infty \frac{(35k-29k+6)3^k}{k(3k-1)(3k-2)\binom{4k}{k}}=\pi\sqrt3,
\label{Au}
\end{equation}
and Sun \cite{series2} proved that
\begin{equation}
\sum_{k=1}^\infty \frac{415k^3-343k+62}{k(3k-1)(3k-2)(-8)^k\binom{4k}{k}}=-3\log2.
\label{Sun}   
\end{equation}
Motivated by these results, we use new methods to determine the values of some of these series.
\begin{theorem}
    We have the identity
    \[\sum_{k=1}^\infty  \frac{(13k^2-15k+2)9^{k-1}}{k(3k-1)(3k-2)\binom{4k}{k}}=\frac{\pi\sqrt3}{3}.\]
\end{theorem}
This first theorem proves Sun's conjecture \cite{series2}*{(4.20)}. 

Let $H_n:=\sum_{k=1}^n 1/k \quad (n\in\N)$ be the $n$-th harmonic number. Define
\[K:=L(2,\chi_{-3})=\sum_{n=0}^\infty\left(\frac{1}{(3n+1)^2}-\frac{1}{(3n+2)^2}\right)=\frac{4\sqrt3}{9}\text{Cl}_2(\pi/3).\]
Here, $\text{Cl}_2$ denotes the Clausen function defined on $\C$ by:
\[\text{Cl}_2(z)=\sum_{k=1}^\infty\frac{\sin(kz)}{k^2}.\]
Calegari, Dimitrov, and Tang \cite{cdt} proved that the constant $K$ is irrational. 

The second theorem of this article is as follows:
\begin{theorem}
    Let $P(X)=95X^2-84X+16$. Then,
    \[\sum_{k=1}^\infty  \frac{P(k)(H_{4k-1}-H_{k-1})-25k+12}{k(3k-1)(3k-2)\binom{4k}{k}}\left(\frac98\right)^{k-1}=\frac{\pi\sqrt3}{3}\log{3}+\frac{15}{4}K.\]
\end{theorem}
This result proves Sun's conjecture \cite{series2}*{(4.19)}.

Our proofs primarily rely on the Beta function $B(x,y)$, defined for every $(x,y)\in\R^2$ by:
\[B(x,y):=\frac{\Gamma(x)\Gamma(y)}{\Gamma(x+y)}\]
where $\Gamma(x)=\int_0^{+\infty}t^{x-1}e^{-t}\,dt$ is the usual Gamma function. The Beta function also has an Euler integral representation given by:
\begin{equation}
B(x,y)=\int_0^1t^{x-1}(1-t)^{y-1}\,dt.
\label{beta_int}
\end{equation}
We use standard notation throughout.
%%%%%%%%%%%%%%%%%%%%%%%%%%%%%%%%%%%%%%%%%%%%%%%%%%%%%%%%%%%%%%%%%%%%%%%
\section{Preliminary results}
The proofs of the theorems in this note rely on a few intermediate results. The first lemma provides a decomposition for quadratic polynomials.
\begin{lemma}
    For every real polynomial $P$ of degree at most $2$, there exists a triple $(\alpha,\beta,\gamma)\in\R^3$ such that:
    \[P(X)=\alpha (3X-1)(3X-2)+\beta X(X+1)+\gamma X(4X-1)\]
\end{lemma}
Note that $\alpha$, $\beta$, and $\gamma$ can be readily determined by solving a system of linear equations. We will compute these values as needed. 
\begin{proof}
    Consider the vector space $\R_2[X]$ of polynomials of degree at most $2$. Define the polynomials:
    \[P_1=(3X-1)(3X-2),\quad P_2=X(X+1),\quad P_3=X(4X-1).\]
    It is straightforward to verify that $(a,b,c)=(0,0,0)$ is the unique solution to the equation:
    \[aP_1+bP_2+cP_3=0\]
    Therefore, $P_1$, $P_2$, and $P_3$ are linearly independent in $\R_2[X]$.
    
    Moreover, since $\dim\R_2[X]=3$, the set $(P_1,P_2,P_3)$ forms a basis for $\R_2[X]$. 

    Consequently, every polynomial $Q\in\R_2[X]$ can be uniquely written as a linear combination of $P_1$, $P_2$, and $P_3$.
\end{proof}
Our second lemma introduces a useful transformation used in the proof of Lemma 2.
\begin{lemma}
    For all integers $n,m\ge1$, 
    \begin{equation}
     \frac{H_{n+m-1}-H_{n-1}}{n\binom{n+m-1}{n}}=-\frac{\partial}{\partial x}B(x,y)\Bigg|_{\substack{x=n\\y=m}}.
    \label{lem2}
    \end{equation}
\end{lemma}
\begin{proof}
    Observe that
    \begin{align*}
        \frac{\partial}{\partial x}B(x,y)&=\frac{\partial}{\partial x}\frac{\Gamma(x)\Gamma(y)}{\Gamma(x+y)}\\&=\frac{\Gamma(x)\Gamma(y)}{\Gamma(x+y)}(\psi(x)-\psi(x+y))\\&=B(x,y)(\psi(x)-\psi(x+y))
    \end{align*}
    where $\psi(x)=\frac{d}{dx}\log{\Gamma(x)}$. Evaluating at $(x,y)=(n,m)$ for integers $n,m\ge1$, we obtain:
    \begin{align*}
        \frac{\partial}{\partial x}B(x,y)\Bigg|_{\substack{x=n\\y=m}}&=B(n,m)(\psi(n)-\psi(n+m))\\&=\frac{1}{n\binom{n+m-1}{n}}(H_{n-1}-H_{n+m-1})
    \end{align*}
    since $\psi(n)=H_{n-1}-\gamma$, where $\gamma$ is the Euler-Mascheroni constant. This proves the lemma.
\end{proof}
The third lemma will be useful at the end of the proof of Theorem 2.
\begin{lemma}
    We have the following identity:
    \begin{equation}
    \int_0^1\frac{\log(1-t)}{3t^2-6t+4}\,dt=-\frac{\pi\sqrt{3}}{18}\log3-\frac{5}{8}K.
    \label{lem3}
    \end{equation}
\end{lemma}
\begin{proof}
    Let 
    \[I=\int_0^1\frac{\log(1-t)}{3t^2-6t+4}\,dt.\]
    Make the change of variable $t=1-x$. Then,
    \[I=\int_0^1\frac{\log(x)}{3x^2+1}\,dx.\]
    Now, making the change of variable $x\sqrt3=\tan\theta$, we obtain:
    \begin{align}
        I&=\frac{\sqrt3}{3}\int_0^{\pi/3}\log\left(\frac{\tan\theta}{\sqrt3}\right)\,d\theta\nonumber\\
        &=-\frac{\pi\sqrt{3}}{18}\log3+\frac{\sqrt3}{3}\left(\int_0^{\pi/3}\log\left(\sin\theta\right)\,d\theta-\int_0^{\pi/3}\log\left(\cos\theta\right)\,d\theta\right)\nonumber\\
        &:=-\frac{\pi\sqrt{3}}{18}\log3+\frac{\sqrt3}3(J_1-J_2)
        \label{I=J1-J2}
    \end{align}
    where $J_1$ and $J_2$ are the respective integrals in the second line. Recall the integral definition of the Clausen function:
    \[\text{Cl}_2(z)=-\int_0^z\log\left|2\sin\frac{\theta}{2}\right|\,d\theta.\]
    Making the change of variable $\theta=u/2$ in the expression for $J_1$, it is straightforward to verify that:
    \begin{equation}
        J_1=-\frac12\text{Cl}_2(2\pi/3)-\frac\pi3\log2.
    \label{J1}
    \end{equation}
    Moreover, making the change of variable $u=\pi/2-\theta$ in the expression for $J_2$, we find:
    \[J_2=\int_{\pi/6}^{\pi/2}\log{(\sin u)}\,du.\]
    This can be evaluated in the same manner as $J_1$. Thus, after some simplification,
    \begin{equation}
    J_2=\frac12\text{Cl}_2(\pi/3)-\frac\pi3\log2.
    \label{J2}
    \end{equation}
    Since \[\text{Cl}_2(\pi/3)=\frac32\text{Cl}_2(2\pi/3)=\frac{3\sqrt3}{4}K,\]
    substituting \eqref{J1} and \eqref{J2} into \eqref{I=J1-J2}, we obtain:
    \begin{align*}
        I&=-\frac{\pi\sqrt{3}}{18}\log3-\frac{\sqrt3}3\cdot\frac{5\sqrt3}{8}K\\&=-\frac{\pi\sqrt{3}}{18}\log3-\frac{5}{8}K.
    \end{align*}
    This proves Lemma 3.
\end{proof}

%%%%%%%%%%%%%%%%%%%%%%%%%%%%%%%%%%%%%%%%%%%%%%%%%%%%%%%%%%%%%%%%%%
\section{Proof of Theorem 1}
We can now prove the first theorem of this article.
\begin{proof}[Proof of Theorem 1]
    Let
    \[S:=\sum_{k=1}^\infty \frac{(13k^2-15k+2)9^{k-1}}{k(3k-1)(3k-2)\binom{4k}{k}}.\]
    The triple $(\alpha,\beta,\gamma)$ associated with the polynomial $13X^2-15X+2$ in the decomposition of Lemma 1 is $(1,-4,2)$. Consequently,
    \[S=\sum_{k=1}^\infty \frac{9^{k-1}}{k\binom{4k}{k}}-4\sum_{k=1}^\infty \frac{(k+1)9^{k-1}}{(3k-1)(3k-2)\binom{4k}{k}}+2\sum_{k=1}^\infty \frac{(4k-1)9^{k-1}}{(3k-1)(3k-2)\binom{4k}{k}}.\]
    The terms of each sum can thus be expressed using the Beta function $B(x,y)$:
    \[12S=\sum_{k=1}^\infty B(3k,k)9^k-4\sum_{k=1}^\infty B(3k-2,k+2)9^k+2\sum_{k=1}^\infty B(3k-2,k+1)9^k.\]
    Using the integral representation of the Beta function, we obtain:
    \begin{multline*}
        12S=\int_0^1\sum_{k=1}^\infty t^{3k-1}(1-t)^{k-1}9^k\,dt-4\int_0^1\sum_{k=1}^\infty t^{3(k-1)}(1-t)^{k+1}9^k\,dt\\+2\int_0^1\sum_{k=1}^\infty t^{3(k-1)}(1-t)^k9^k\,dt.
    \end{multline*}
    Recognizing geometric series, we have:
    \begin{align}
        \frac43S&=\int_0^1\frac{t^2}{1-9t^3(1-t)}\,dt-4\int_0^1\frac{(1-t)^2}{1-9t^3(1-t)}\,dt+2\int_0^1\frac{1-t}{1-9t^3(1-t)}\,dt\nonumber
        \\&=\int_0^1\frac{-3t^2+6t-2}{9t^4-9t^3+1}\,dt.
        \label{S=Int}
    \end{align}
    Let $u$ and $v$ be two polynomials defined on $\R$ by:
    \[u(x)=3x^2-\frac32x-\frac12\quad\text{and}\quad v(x)=\frac{\sqrt3}{2}(x-1).\]
    One can verify that:
    \begin{equation}
    \frac{-3x^2+6x-2}{9x^4-9x^3+1}=-\frac{2\sqrt3}{3}\cdot\frac{d}{dx}\arctan{\frac{u(x)}{v(x)}}.
    \label{d/dx atan u/v}
    \end{equation}
    \[
    \left|
    \hspace{0.2cm}
    \begin{minipage}{\textwidth}
    \textbf{Remark.} We determined the polynomials $u$ and $v$ using \texttt{Mathematica}. Note that we need to find $u$ and $v$ such that $u'v-uv'=C(3X^2-6X+2)$ and $u^2+v^2=9X^4-9X^3+2$, where $C$ is a constant. Consequently, one of the two polynomials must have degree $2$ and the other degree $1$ (here, $u$ and $v$, respectively). Moreover, by coefficient comparison, we may set the coefficient of $X^2$ in the polynomial $u$ equal to $3$. This reduces the number of unknowns in the system.
    \end{minipage}
    \hspace{0.2cm}
    \right|
    \]\\
 
    Thus, by substituting \eqref{d/dx atan u/v} into \eqref{S=Int}, we obtain the desired result:
    \begin{align*}
        S&=\lim_{x\to1^-}-\frac34\times\frac{2\sqrt3}{3}\left(\arctan\frac{u(x)}{v(x)}-\arctan\frac{u(0)}{v(0)}\right)\\
        &=-\frac{\sqrt{3}}{2}\left(-\frac{\pi}2-\frac\pi6\right)\\
        &=\frac{\pi\sqrt3}3.
    \end{align*}
    This concludes the proof.
\end{proof}
%%%%%%%%%%%%%%%%%%%%%%%%%%%%%%%%%%%%%%%%%%%%%%%%%%%%%%%%%%%%%%%%%%%%%%%%%%%%%%%%%%%
\section{Proof of Theorem 2}
We can now prove our second theorem. This proof partially follows the arguments used in the proof of Theorem 1.
\begin{proof}[Proof of Theorem 2]
    Let $P(X)=95X^2-84X+16$ be a polynomial. We have
    \begin{equation}
    S:=\sum_{k=1}^\infty  \frac{P(k)(H_{4k-1}-H_{k-1})-25k+12}{k(3k-1)(3k-2)\binom{4k}{k}}\left(\frac98\right)^{k-1}:=\frac89(S_1+S_2)
    \label{S}
    \end{equation}
    where
    \[S_1:=\sum_{k=1}^\infty  \frac{P(k)(H_{4k-1}-H_{k-1})}{k(3k-1)(3k-2)\binom{4k}{k}}\left(\frac98\right)^k\]
    and
    \[S_2:=\sum_{k=1}^\infty  \frac{-25k+12}{k(3k-1)(3k-2)\binom{4k}{k}}\left(\frac98\right)^k.\]
    The triple $(\alpha,\beta,\gamma)$ associated with the polynomial $P$ in the decomposition of Lemma 1 is $(8,-5,7)$. Thus,
    \begin{multline*}
        S_1=8\sum_{k=1}^\infty \frac{H_{4k-1}-H_{k-1}}{k\binom{4k}{k}}\left(\frac98\right)^k-5\sum_{k=1}^\infty\frac{(k+1)(H_{4k-1}-H_{k+1})}{(3k-1)(3k-2)\binom{4k}{k}}\left(\frac98\right)^k\\-5\sum_{k=1}^\infty\frac{(k+1)\left(\frac1k+\frac1{k+1}\right)}{(3k-1)(3k-2)\binom{4k}{k}}\left(\frac98\right)^k+7\sum_{k=1}^\infty\frac{(4k-1)(H_{4k-2}-H_{k})}{(3k-1)(3k-2)\binom{4k}{k}}\left(\frac98\right)^k\\+7\sum_{k=1}^\infty\frac{(4k-1)\left(\frac1{4k-1}+\frac1k\right)}{(3k-1)(3k-2)\binom{4k}{k}}\left(\frac98\right)^k.
    \end{multline*}
    Set
    \[S_3:=\sum_{k=1}^\infty\frac{(k+1)\left(\frac1k+\frac1{k+1}\right)}{(3k-1)(3k-2)\binom{4k}{k}}\left(\frac98\right)^k\]
    and
    \[S_4:=\sum_{k=1}^\infty\frac{(4k-1)\left(\frac1{4k-1}+\frac1k\right)}{(3k-1)(3k-2)\binom{4k}{k}}\left(\frac98\right)^k.\]
    Note that
    \[5S_3-7S_4=\sum_{k=1}^\infty \frac{-25k+12}{k(3k-1)(3k-2)\binom{4k}{k}}\left(\frac98\right)^k=S_2.\]
    By the transformation \eqref{lem2} from Lemma 2, we have:
    \begin{multline*}
        -\frac43S_1=8\sum_{k=1}^\infty \left(\frac98\right)^k\frac{\partial}{\partial y}B(x,y)\Bigg|_{\substack{x=3k\\y=k}}-5\sum_{k=1}^\infty\left(\frac98\right)^k\frac{\partial}{\partial y}B(x,y)\Bigg|_{\substack{x=3k-2\\y=k+2}}\\+7\sum_{k=1}^\infty\left(\frac98\right)^k \frac{\partial}{\partial y}B(x,y)\Bigg|_{\substack{x=3k-2\\y=k+1}}+\frac{4}{3}S_2.
    \end{multline*}
    Differentiating the integral representation \eqref{beta_int} of the Beta function with respect to $y$, we obtain:
    \begin{align*}
        -\frac43(S_1+S_2)=&8\int_0^1\sum_{k=1}^\infty t^{3k-1}(1-t)^{k-1}\left(\frac98\right)^k\log{(1-t)}\,dt\\&-5\int_0^1\sum_{k=1}^\infty t^{3(k-1)}(1-t)^{k+1}\left(\frac98\right)^k\log(1-t)\,dt\\&+7\int_0^1\sum_{k=1}^\infty t^{3(k-1)}(1-t)^k\left(\frac98\right)^k\log(1-t)\,dt
        \\=&\frac98\Bigg(8\int_0^1\frac{t^2}{1-\frac98t^3(1-t)}\log(1-t)\,dt\\&-5\int_0^1\frac{(1-t)^2}{1-\frac98t^3(1-t)}\log(1-t)\,dt\\&+7\int_0^1\frac{1-t}{1-\frac98t^3(1-t)}\log(1-t)\,dt\Bigg)\\=&\int_0^1\frac{9(3t^2+3t+2)}{9t^4-9t^3+8}\log(1-t)\,dt.
    \end{align*}
    Notice that $9X^4-9X^3+8=(3X^2+3X+2)(3X^2-6X+4)$. Then,
    \[-\frac4{27}(S_1+S_2)=\int_0^1\frac{\log(1-t)}{3t^2-6t+4}\,dt.\]
    Thus, the identity \eqref{lem3} of Lemma 3 yields:
    \[S_1-S_2=\frac{3\pi\sqrt{3}}{8}\log3+\frac{135}{32}K.\]
    Hence,
    \begin{equation}
    S_1+S_2=\frac{3\pi\sqrt{3}}{8}\log3+\frac{135}{32}K.
    \label{S1}
    \end{equation}
    Then, by substituting \eqref{S1} into \eqref{S}, we arrive at the expected result:
    \begin{align*}
        S&=\frac89(S_1+S_2)\\
        &=\frac89\left(\frac{3\pi\sqrt{3}}{8}\log3+\frac{135}{32}K\right)\\
        &=\frac{\pi\sqrt3}3\log{3}+\frac{15}{4}K.
    \end{align*}
    This proves Theorem 2.
    \end{proof}
\section{Concluding remarks}
Our method also provides a new proof of the identities \eqref{Au} and \eqref{Sun}. Sun also conjectured the following formula. 
\begin{conj}[Sun \cite{series2}*{(4.25)}]
Let $P(X)=415X^2-343X+62$ and $Q(X)=581X^2-229X-6$. Then,
\[\sum_{k=1}^\infty \frac{P(k)-Q(k)/(4k)}{k(3k-1)(3k-2)(-8)^k\binom{4k}{k}}=-\frac{\pi^2}{4}.\]
\end{conj}
Apparently, our method does not allow us to prove this identity. Moreover, we cannot prove the similar formulas involving $H_{2k-1}$ (see \cite{series2}*{(4.21)-(4.24)}).
\bibliographystyle{amsplain}
\bibliography{bibly}
\end{document}